\documentclass[10pt]{article}
\usepackage{icmlko}

\icmlmeta{IP 파운데이션 월드모델}{289}{아무도 안 주웠다}{누출이 놓여 있었는데 모형이 안 쓴다 --- 표지를 지워도 0.0033 밖에 안 잃는다}

\begin{document}

\twocolumn[
\icmlheader

\begin{abstract}
노트 288이 넓힌 팝업 판에서 출처 표지를 찾았다 --- 유보에서 관측 표시자
하나로 $|\rho|{=}0.4157$ 이고 같은 판의 실제 팝업 점수 0.3582 보다 크다.
``표지가 모형을 이긴다''고 적고 노트 127 · 133의 채택 근거를 다시 보라고
남겼다. \textbf{쟀다.} 마스크를 상수로 만들어 표지를 지우니 팝업이
\textbf{$-$0.0033} 이다(0.3582 $\to$ 0.3549). \textbf{쓸 수 있는 0.4157 중에
실제로 쓰인 것이 0.0033 --- 126배 차이다.} \textbf{누출이 놓여 있었는데
아무도 안 주웠다.} 그래서 노트 127 · 133의 이득은 표지가 아니다. 그리고
같은 실험에서 시장 행을 빼니 팝업이 0.3582 $\to$ 0.3980 으로 크게 올랐는데,
\textbf{그건 채점 대상이 116행에서 65행으로 바뀐 것}이라 그대로 못 읽는다.
\textbf{같은 내부 유보 65행에 고정}하고 시장 행을 학습에서만 빼니
\textbf{$+$0.0028} --- 학습행 56개가 사라지는데 아무 일도 안 난다. 분해가
닫힌다: 표지 기여 $-$0.0033 · 학습 기여 $+$0.0028 · \textbf{나머지 전부가
채점 대상 차이}다. 결론 --- \textbf{넓힌 판의 시장 107행은 학습에 아무
기여도 안 하면서 채점만 어렵게 만든다.} 노트 283 · 284가 ``붙이지 말고 따로
열라''고 한 결정이 \textbf{두 번째 독립된 이유로} 옳았다. 첫 번째는 표지가
생긴다는 것이었고, 이번 것은 \textbf{붙여도 얻는 게 없다}는 것이다.
\end{abstract}
]

\section{쓸 수 있는 것과 실제로 쓴 것}

노트 288이 넓힌 팝업 판(189행 $=$ 내부 82 $+$ 시장 107)에서 손 축 마스크가
내부 98.8\% 대 시장 0.0\% 이고, 그 표시자 하나로 유보 $|\rho|{=}0.4157$ 이
나오는 것을 보였다. 실제 팝업 점수가 0.3582 이니 \textbf{표지가 모형보다
잘 맞힌다.}

\textbf{그런데 모형이 그걸 쓰고 있나?} 표지를 지워 봤다 --- 손 축 다섯의
관측 표시자를 전부 1 로 놓으면 값은 그대로(시장 행은 중립 0.5)인데 표지가
사라진다.

\begin{figure*}[t]
\centering
\includegraphics[width=\textwidth]{figs/picked.pdf}
\caption{왼쪽 --- 표지 하나로 낼 수 있는 값과 표지를 지웠을 때 실제로 잃는
값. 오른쪽 --- 네 판의 팝업 $\rho$. 위 둘은 유보 116행, 아래 둘은 내부 유보
65행에 고정한 것.}
\end{figure*}

\begin{center}\small
\begin{tabular}{lrrl}
\hline
& 팝업 $\rho$ & 판 & 출처 가드 \\
\hline
그대로 & 0.3582 & 0.4076 & \textbf{실패} \\
\textbf{표지 지움} & \textbf{0.3549} & 0.4098 & 통과 \\
\hline
\multicolumn{4}{l}{차 \textbf{$-$0.0033} (판은 오히려 $+$0.0022)} \\
\hline
\end{tabular}
\end{center}

\textbf{0.4157 을 낼 수 있는 열을 지웠는데 0.0033 을 잃는다. 126배다.}

\paragraph{왜 안 썼나} F18은 열한 도메인을 한 표로 놓고 적합하는 배깅
부스팅이다. 팝업은 유보 2,700여 행 중 116행이고, 그 표지는 스물세 열 중
다섯의 지표다. \textbf{나무가 팝업 하나를 위해 특화할 이유가 없다.} 노트
259가 잰 것의 다른 얼굴이다 --- 분기는 열한 도메인이 함께 내려가는 길이라
한 도메인에만 쓸모 있는 열은 뿌리 근처에 못 온다.

\paragraph{그러니 노트 127 · 133은 무사하다} 그 노트들이 넓힌 판의 팝업
이득($+$0.3688 $\to$ $+$0.4495)을 ``팝업이 학습 풀에 들어가서''로 설명했고,
노트 288이 ``출처 표지도 같은 크기를 낸다''는 대안을 세웠다. \textbf{그
대안이 지금 판에서 0.0033 이다.} 다른 판 · 다른 축이었으므로 완전한 면죄는
아니지만, \textbf{크기가 안 맞는다.}

\section{그리고 첫 수를 그대로 읽으면 안 됐다}

같은 실험에서 시장 행을 통째로 빼 봤더니 팝업이 \textbf{0.3582 $\to$
0.3980} 이었다. $+$0.0398 이면 큰 수다.

\textbf{그대로 읽으면 안 된다.} 시장 행을 빼면 \textbf{채점 대상도} 116행에서
65행으로 바뀐다. 두 수는 다른 문제의 답이다.

그래서 \textbf{내부 유보 65행에 고정}하고 학습에서만 시장 행을 뺐다.

\begin{center}\small
\begin{tabular}{lrrr}
\hline
& 학습 & 유보 & 팝업 $\rho$ \\
\hline
시장 학습 포함 & 73 & 65 & $+$0.3952 \\
시장 학습 뺌 & \textbf{17} & 65 & $+$0.3980 \\
\hline
\multicolumn{4}{l}{차 \textbf{$+$0.0028}} \\
\hline
\end{tabular}
\end{center}

\textbf{학습행 56개가 사라지는데 아무 일도 안 난다.} 노트 133이 넓힌 판을
만든 목적이 바로 그 행들인데, 지금 판에서는 기여가 0 이다.

\paragraph{분해가 닫힌다}

\begin{center}\small
\begin{tabular}{lr}
\hline
표지가 하는 일 & $-$0.0033 \\
시장 학습행이 하는 일 & $+$0.0028 \\
\hline
\textbf{나머지 전부} & \textbf{채점 대상 차이} \\
\hline
\end{tabular}
\end{center}

0.3582(유보 116, 시장 51 포함)와 0.3952(유보 65, 내부만)의 차이는
\textbf{모형이 아니라 무엇을 채점했나}다. 시장 유보 행이 순위 매기기 더
어렵다.

\section{그래서 붙이지 말라는 결론이 두 번째 이유를 얻는다}

노트 283이 시장 층을 팝업에 붙이지 말라고 한 근거는 \textbf{표지가 생긴다}는
것이었다. 이번 노트가 두 번째를 더한다.

\begin{center}\small
\begin{tabular}{ll}
\hline
붙이면 안 되는 이유 & \\
\hline
① 표지가 생긴다(노트 283 · 288) & 지금 모형은 안 쓰지만 놓여 있다 \\
\textbf{② 붙여도 얻는 게 없다} & \textbf{학습 기여 $+$0.0028} \\
③ 채점만 어려워진다 & 유보에 어려운 행 51개가 섞인다 \\
\hline
\end{tabular}
\end{center}

\textbf{①은 잠재적 위험이고 ②③은 실제 비용이다.} 노트 284가 열한 번째
도메인으로 따로 연 결정은 셋 다 피한다 --- 시장 행은 제 도메인에서
$\rho{=}0.197$ 로 정직하게 채점되고, 팝업은 제 65행으로 채점된다.

\paragraph{넓힌 판을 어떻게 하나} \texttt{wide*} 축 세트에서 시장 107행을
빼는 것이 맞다. 얻는 게 없고 표지를 만들고 채점을 흐린다. 다만 그건 노트
133이 만든 판을 고치는 일이라 \textbf{이 노트에서 바로 안 한다} --- 넓힌
판을 쓴 옛 노트들의 수가 바뀌므로 노트 284의 규약대로 기준선을 다시 놓는
일이 따라온다.

\paragraph{남은 것}
\begin{center}\small
\begin{tabular}{ll}
\hline
갈래 & 왜 \\
\hline
\texttt{wide*} 정리 & 시장 107행 제거 · 기준선 재설정 \\
다른 정식화로 & F6 · F9 는 표지를 쓸까(청력이 높다) \\
표지를 일부러 준다면 & 능형은 0.4157 을 주울지도 \\
노트 133 재현 & 옛 축 세트로 그때 수를 다시 \\
\hline
\end{tabular}
\end{center}

\paragraph{규약} \textbf{누출이 있다는 것과 누출이 쓰인다는 것은 다르다.}
가드는 있는 것을 막아야 하지만(다른 정식화가 주울 수 있다), \textbf{점수를
설명할 때는 실제로 쓰였는지 재야 한다.} 그리고 \textbf{행을 빼는 실험은
채점 대상이 같은지 먼저 본다} --- $+$0.0398 이 $+$0.0028 이 됐다.

\end{document}
